\documentclass{article}
\usepackage[ngerman]{babel}
\usepackage{amsmath}

\begin{document}

\section*{Beschreibung der Skizze}

Die dargestellte Skizze zeigt den Grundriss eines rechteckigen Gebäudes mit einer Gesamtbreite von $6{,}90\,\text{m}$ und einer Gesamtlänge von $10{,}00\,\text{m}$.
Die Außenwände sind als dicke, schwarze Linien dargestellt.

Der Innenraum ist in mehrere Bereiche unterteilt:
\begin{itemize}
   \item Auf der linken Seite befindet sich ein großer Raum mit einer Breite von $3{,}15\,\text{m}$ und einer Länge von $9{,}00\,\text{m}$.
   \item Auf der rechten Seite sind zwei kleinere Räume angeordnet, jeweils mit einer Breite von $2{,}50\,\text{m}$.
   \item Der obere rechte Raum hat eine Länge von $6{,}00\,\text{m}$, während der untere rechte Raum eine Länge von $2{,}75\,\text{m}$ aufweist. 
\end{itemize}

Zwischen den Räumen verlaufen Innenwände, die ebenfalls als dicke Linien eingezeichnet sind. Türen sind durch Kreisbögen dargestellt, die die Öffnungsrichtung anzeigen.

Zusätzliche Markierungen:
\begin{itemize}
   \item Kreise mit dem Buchstaben ``S'' kennzeichnen Schalter.
   \item Kreise mit dem Buchstaben ``X'' markieren Lampen.
\end{itemize}

Weitere Maße innerhalb der Skizze:
\begin{itemize}
   \item Es folgt eine Bescheibung der Wandabstände, der Türen und Fenstern. 
   \item Alle Fenster haben einen Abnstand vom Boden von $1{,}10\,\text{m}$. Deren Höhe beträgt $1{,}30\,\text{m}$. 
   \item Die restliche Höhe von $0{,}40\,\text{m}$ ist Wandmaterial.
   \item Im großen Raum an der linken Seite ist der Abstand von der Tür zur Außenwand $0{,}50\,\text{m}$.
   \item Nach der Tür mit $1{,}00\,\text{m}$ länge ist nochmals eine Wand mit $1{,}50\,\text{m}$.
   \item Danach kommt ein Fenster mit einer länge von $2{,}00\,\text{m}$. Der restliche Raum hat eine Wand mit $4{,}00\,\text{m}$ länge. 
   \item Im oberen rechten Raum dessen linken Wand hat einen Abstand zur Tür von $1{,}00\,\text{m}$ danach folgt eine Tür mit $1{,}00\,\text{m}$ länge.
   \item Die restliche Seite hat eine länge von $4{,}00\,\text{m}$.
   \item Gegenüberliegend davon ist eine Außenwand mit $2{,}00\,\text{m}$ drauf folgend ein Fentser mit der länge von $2{,}00\,\text{m}$.
   \item Danach nochmals eine Wand mit $2{,}00\,\text{m}$. 
   \item Im unteren rechten Raum dessen linken Wand hat einen Abstand zur Tür von $0{,}75\,\text{m}$.
   \item Die Tür wiederum hat eine Länge von $1{,}00\,\text{m}$ danach folgt eine Wand mit $1{,}00\,\text{m}$.
   \item In diesem Raum ist das Fentser an der unteren Außenwand plaziert.
   \item Mit einem Abstand $1{,}00\,\text{m}$ danach folgt ein Fenster mit der Länge von $1{,}00\,\text{m}$ danach eine Wand mit $0{,}50\,\text{m}$.
   \item Die Höhe des Gabäudes ist $2{,}80\,\text{m}$.
\end{itemize}

Die Skizze stellt somit einen strukturierten Grundriss mit klar definierten Raumaufteilungen und elektrischen Installationen dar.

\end{document}